\documentclass[11pt]{article}
\usepackage[margin=1in]{geometry}
\usepackage[T1]{fontenc}
\usepackage{lmodern,microtype}
\usepackage{amsmath,amssymb,amsthm,mathtools,mathrsfs}
\clubpenalty=10000
\widowpenalty=10000
\displaywidowpenalty=10000
\usepackage{hyperref}
\hypersetup{pdfauthor={},pdftitle={Positive physical Hadamard states for critical Einstein--scalar multiplets on compact Cauchy surfaces},colorlinks=true,linkcolor=blue,citecolor=blue,urlcolor=blue}
\newtheorem{theorem}{Theorem}
\newtheorem{proposition}[theorem]{Proposition}
\newtheorem{lemma}[theorem]{Lemma}
\newtheorem{corollary}[theorem]{Corollary}
\theoremstyle{definition}
\newtheorem{remark}[theorem]{Remark}
\DeclareMathOperator{\Ran}{Ran}
\DeclareMathOperator{\Ker}{Ker}
\DeclareMathOperator{\WF}{WF}
\DeclareMathOperator{\tr}{tr}
\newcommand{\cH}{\mathcal H}
\newcommand{\cV}{\mathcal V}
\newcommand{\cN}{\mathcal N}
\newcommand{\id}{\mathbf 1}
\title{Positive physical Hadamard states for critical Einstein--scalar multiplets on compact Cauchy surfaces}
\author{}
\date{29 October 2026}
\begin{document}
\maketitle
\begin{abstract}
We construct positive, gauge-invariant quasifree Hadamard states for linearized Einstein gravity coupled to finitely many canonical real scalar fields.
The background is a smooth four-dimensional globally hyperbolic solution with compact Cauchy surface, the scalar target metric is fixed and positive definite, and the potential gradient vanishes along the background image.
The scalar fields may be nonconstant, and the potential Hessian need not vanish or be positive.
We derive the finite-multiplet Hessian and gauge complex in one common normalization, establish compatible microlocal projections and an exact positive metric gauge slice, and treat the full Fredholm obstruction by a smooth finite-rank correction formed with the physical charge.
The resulting covariances descend to the physical quotient and lift to spacetime with the Hadamard wavefront relation.
A closed Friedmann--Lema\^itre--Robertson--Walker example has nonconstant scalar fields, nonzero metric--scalar mixing and, after adjoining a spectator, a nonzero potential Hessian.
\end{abstract}

\section{Introduction and theorem}
Positive state construction for a gauge theory requires compatibility between the field equations, gauge reduction and microlocal regularity.
For coupled metric and matter perturbations, formal gauge invariance and normal hyperbolicity alone do not supply positivity on the physical observables.
The indefinite gravitational charge must be controlled on a genuine gauge slice, including the smooth finite-dimensional obstructions to complete gauge fixing.
This paper proves that the compact-Cauchy construction for linearized gravity extends to a finite canonical scalar multiplet under a critical-potential hypothesis along the background image.

The covariant linearized Einstein--Klein--Gordon framework is established in Hack~\cite{Hack2014EKG}.
The one-real-scalar complex in the normalization used below is displayed by Schmid~\cite[Section 2.2.5]{Schmid2026Thesis}.
The microlocal construction, complete gauge fixing and finite-obstruction architecture are inherited from G\'erard's compact-Cauchy gravity theorem~\cite{Gerard2024Compact}, using pseudodifferential Hadamard methods developed in~\cite{GerardWrochna2014Calculus,GerardWrochna2015YangMills,GerardOulghaziWrochna2017Geometry}.
The coupled compatibility proofs below replace the vacuum-Einstein identities that cannot be transferred directly.
The asymptotically flat gravity result~\cite{GerardMolodykWrochna2026Flat} concerns a different background class.
No claim of priority over the complete coupled-gauge literature is needed here.

\begin{theorem}[Positive physical Hadamard states]\label{thm:main}
Let $(M,g,\Phi)$ be a smooth, time-oriented, four-dimensional globally hyperbolic on-shell solution of canonical Einstein gravity coupled to $N<\infty$ real scalar fields, with a compact smooth Cauchy surface $\Sigma$.
The scalar target metric is a fixed positive-definite matrix.
Suppose that the smooth potential satisfies
\[
V_{,A}(\Phi(x))=0\qquad(x\in M,\ A=1,\ldots,N).
\]
For the linearized gauge complex constructed below, put
\[
\cV_{\mathrm{phys}}=\Ker K_\Sigma^\dagger/\Ran K_\Sigma
\]
on smooth Cauchy data, with its induced physical charge.
There exist positive Hermitian covariances $\lambda_\Sigma^\pm$ on this quotient satisfying the canonical commutation relations.
Their spacetime lifts define a gauge-invariant quasifree Hadamard state for the real linearized coupled field.
\end{theorem}

Here ``positive scalar'' refers to the kinetic target metric.
The matrix $V_{,AB}$ may be indefinite or singular, subject to the identities forced by the hypotheses.
Differentiating criticality along the background yields $V_{,AB}\,d\Phi^B=0$.
Thus a varying scalar background lies along flat directions of the critical set; a transverse mass matrix remains allowed.
No constant-background or rank assumption on $d\Phi$ is imposed.

The theorem is an existence result for the linearized canonical coupled system.
It does not select a state or assert uniqueness, purity, nonlinear quantization or a dynamical stability conclusion.
The proof uses compactness for the Fredholm and finite-obstruction steps and retains criticality in the trace-reversal argument.
The result is not asserted for a field-dependent target metric, a general noncritical potential or an arbitrary noncompact Cauchy surface.

\paragraph{Analytic inputs and proof organization.}
We use the well-posed Cauchy problem, retarded and advanced Green operators and propagation of singularities for smooth normally hyperbolic operators on finite-rank bundles.
The pseudodifferential input is the Hermitian-bundle Riccati construction of~\cite[Sections 3.1--3.2]{Gerard2024Compact}, based on~\cite[Section 5.2.7, Proposition 5.2]{GerardMurroWrochna2025Wick}: scalar positive principal root, exact graph projections, smoothing frequency leakage and elliptic regularity on compact $\Sigma$.
The proof derives the coupled Hessian, exact charges, symmetry, gauge-slice positivity, Fredholm representative and finite-rank repair explicitly.
Consequently the new compatibility argument is not replaced by an appeal to a generic parametrix.

\section{Action, Hessian and common normalization}

For clarity about physical units, begin more generally with

\[
 S_{\rm phys}=\int_M\left[\frac{R-2\Lambda}{2\kappa}
 -\frac12 G_{AB}\nabla\phi^A\cdot\nabla\phi^B-V_{\rm phys}(\phi)\right]d\mathrm{vol}_g,
 \qquad \kappa>0,
\]

where $G$ is a fixed positive-definite matrix. The linear change of scalar coordinates $\Phi=\sqrt\kappa\,G^{1/2}\phi$ and potential $V(\Phi)=\kappa V_{\rm phys}(G^{-1/2}\Phi/\sqrt\kappa)$ gives $S_{\rm phys}=S/\kappa$ with $S$ in (1). Thus the physical Hessian is $1/\kappa$ times the Hessian below. This positive common overall factor is distinct from the relative metric/scalar normalization in (8); it is restored by the corresponding positive normalization of the CCR charge. ``Positive scalar multiplet'' refers to the target kinetic metric $G$, not to the sign of $V_{AB}$. The potential Hessian may be indefinite or singular, subject only to smoothness and the identities forced by a critical on-shell background.

Let the signature be $(-+++)$, let $N<\infty$, and absorb the positive constant gravitational coupling and a constant positive scalar target metric into canonical scalar coordinates. Start with

\[
 S[g,\Phi]=\int_M\left[\frac12(R-2\Lambda)
 -\frac12\sum_A g^{ab}\partial_a\Phi^A\partial_b\Phi^A-V(\Phi)\right]d\mathrm{vol}_g.
 \tag{1}
\]

Write $a_A=d\Phi^A$, $H_A=\nabla^2\Phi^A$, $S_{ab}=\sum_Aa_{Aa}a_{Ab}$, and $c=\Lambda+V(\Phi)$. The background equations are

\[
 E_{ab}=\mathrm{Ein}_{ab}+\Lambda g_{ab}
 -S_{ab}+\left(\tfrac12\sum_A|a_A|_g^2+V\right)g_{ab}=0,
 \qquad E_A=-\Box\Phi^A+V_A=0.
 \tag{2}
\]

Consequently $\mathrm{Ric}-S=cg$. Define

\[
 (d_s\xi)_{ab}=\nabla_{(a}\xi_{b)},\qquad
 (\delta_s h)_a=-2\nabla^b h_{ba},\qquad
 Ih=h-\tfrac12g\operatorname{tr}_g h.
 \tag{3}
\]

Thus $d_s^*=\delta_s$ when the one-form pairing is the ordinary Lorentzian contraction and the symmetric two-tensor pairing is

\[
 \langle h,k\rangle_{S_2}=2\overline h_{ab}k^{ab}.
 \tag{4}
\]

The curvature endomorphism convention is $\mathcal R(h)_{ab}=-R_{acbd}h^{cd}$ with the usual Riemann convention for which $g^{cd}R_{acbd}=\mathrm{Ric}_{ab}$. In particular $\mathcal Rg=-\mathrm{Ric}$ and $\operatorname{tr}\mathcal Rh=-\mathrm{Ric}:h$. Set

\[
 (S\times_s h)_{ab}=S_{(a}{}^c h_{b)c},\qquad
 A_g=2\mathcal R+2S\times_s .
 \tag{5}
\]

This specifies the signs independently of names used for curvature operators in different sources.

For original perturbations $(h,\psi)=(\delta g,\delta\Phi)$, let $L=\delta(E_{ab},E_A)$. Direct differentiation of (2) gives

\[
\begin{aligned}
 L_{gg}&=\tfrac12(-\Box-I d_s\delta_s+A_g)I,\\
 L_{gA}\psi^A&=-2I(a_A\odot d\psi^A)+gV_A\psi^A,\\
 (L_{Ag}h)_A&=H_A:h-\tfrac12a_A\cdot\delta_s Ih,\\
 L_{AB}&=-\Box\delta_{AB}+V_{AB}.
\end{aligned}\tag{6}
\]

Repeated scalar indices are summed. Here $\odot$ denotes symmetrization with factor $1/2$. To check the mixed terms without relying on a block ansatz, differentiation of the scalar equation gives

\[
 \delta(-\Box\Phi^A+V_A)
 =-\Box\psi^A+V_{AB}\psi^B
 +h^{ab}H_{Aab}
 +(\nabla^ah_a{}^b-\tfrac12\nabla^b\operatorname{tr}h)a_{Ab}.
 \tag{7}
\]

Differentiation of the scalar stress gives the second line of (6). The standard metric variation

\[
 \delta\mathrm{Ric}_{ab}=-\tfrac12\Box h_{ab}
 -R_{acbd}h^{cd}+\mathrm{Ric}_{(a}{}^ch_{b)c}
 +\nabla_{(a}(\nabla^ch_{b)c}-\tfrac12\nabla_{b)}\operatorname{tr}h)
\]

together with $\mathrm{Ric}-S=cg$ gives the first line of (6). The scalar mass matrix is the full $N\times N$ Hessian $V_{AB}$; no diagonalization or vanishing-mass assumption has been made.

The untransformed Hessian pairing is

\[
 B_0((h,\psi),(k,\chi))=2\overline h:k+4\overline\psi\cdot\chi.
\]

On shell, compactly supported variations satisfy $S''(u,v)=-\tfrac14\int B_0(u,Lv)$. This verifies formal self-adjointness of $L$ for $B_0$. One can also check the mixed adjoint by integration by parts: $L_{gA}^*=4L_{Ag}$ in the $S_2$ and scalar pairings, because the discrepancy is proportional to $-\Box\Phi^A+V_A$.

Now write the original metric perturbation as $h=2I\gamma$ and define

\[
 \mathcal A=\operatorname{diag}(2I,1_N),\qquad
 P=L\mathcal A,\qquad
 B((\gamma,\psi),(k,\chi))
 =4(\overline\gamma:Ik+\overline\psi\cdot\chi).
 \tag{8}
\]

Then $S''(\mathcal Au,\mathcal Av)=-\tfrac14\int B(u,Pv)$, so $P=P^\dagger$ for the physical bundle form $B$. Explicitly,

\[
 P=\begin{pmatrix}
 -\Box-I d_s\delta_s+A_g&-2I(a_B\odot d\,\cdot)+gV_B\\
 2H_A:I(\cdot)-a_A\cdot\delta_s(\cdot)&-\Box\delta_{AB}+V_{AB}
 \end{pmatrix}.
 \tag{9}
\]

In this formula and below the field bundle is $E_2=S^2T^*M_\mathbb C\oplus\mathbb C^N$; the gauge bundle is $E_1=T^*M_\mathbb C$, with the ordinary Lorentzian Hermitian pairing.

\section{The coupled gauge complex}

Set

\[
 K\xi=\frac1{\sqrt2}\begin{pmatrix}I d_s\xi\\a_A\cdot\xi\end{pmatrix},
 \qquad
 K^\dagger(\gamma,\psi)=\sqrt2\left(\delta_s\gamma+2\sum_Aa_A\psi^A\right).
 \tag{10}
\]

The adjoint follows directly from (8), (3) and (4): the metric integration by parts contributes $-2\sqrt2\operatorname{div}\gamma$ and the scalar term contributes $2\sqrt2\sum_Aa_A\psi^A$. Moreover $\mathcal AK\xi$ is the infinitesimal diffeomorphism generated by $X=\xi^\sharp/\sqrt2$: its metric component is $\sqrt2d_s\xi$ and its scalar component is $a_A\cdot\xi/\sqrt2$. Naturality of the equations therefore gives

\[
 PK=0,\qquad K^\dagger P=0.
 \tag{11}
\]

This is an on-shell identity; it has not been assumed for arbitrary prescribed scalar profiles.

Use $\delta_sId_s=-\Box-\mathrm{Ric}_a{}^b$ to obtain

\[
 D_1=K^\dagger K=-\Box-\mathrm{Ric}+2S
 =-\Box-c\,1_{E_1}+S.
 \tag{12}
\]

Here $S$ acts on one-forms as $\xi\mapsto\sum_Aa_A(a_A\cdot\xi)$. The product before cancellation is

\[
 KK^\dagger=
 \begin{pmatrix}
 I d_s\delta_s&2I(a_B\odot d\,\cdot)+2I H_B\\
 a_A\cdot\delta_s&2a_A\cdot a_B
 \end{pmatrix}.
 \tag{13}
\]

Since $2I H_B=2H_B-g\Box\Phi^B=2H_B-gV_B$, all first-order mixed terms cancel in $D_2=P+KK^\dagger$, yielding

\[
 \boxed{D_2=
 \begin{pmatrix}
 -\Box+A_g&2H_B\\
 2H_A:I(\cdot)&-\Box\delta_{AB}+V_{AB}+2a_A\cdot a_B
 \end{pmatrix}.}
 \tag{14}
\]

Equations (10)--(14) extend the one-real-scalar construction of Schmid~\cite[Section 2.2.5]{Schmid2026Thesis} to a finite multiplet. The scalar--scalar $a_A\cdot a_B$ term need not be diagonal. Both $D_i$ are normally hyperbolic; all metric--scalar mixing in $D_2$ is multiplication. Consequently their principal connections are respectively the Levi-Civita connection and the direct sum of the induced tensor connection with the trivial scalar connection. This remains true for nonconstant $\Phi$.

Algebraically,

\[
 D_2K=KD_1,\qquad K^\dagger D_2=D_1K^\dagger.
 \tag{15}
\]

Normal hyperbolicity and uniqueness of retarded/advanced solutions then give $G_2K=KG_1$ and $K^\dagger G_2=G_1K^\dagger$ on compactly supported test sections, for each choice of Green operator and for their causal difference.

\section{Critical potentials and trace reversal}

Put $J=I\oplus1_N$. Since $\nabla I=0$, the rough wave operator commutes with $J$. The curvature-plus-stress block satisfies

\[
 A_gg=-2(\mathrm{Ric}-S)=-2cg,\qquad
 \operatorname{tr}(A_g\gamma)=-2(\mathrm{Ric}-S):\gamma=-2c\operatorname{tr}\gamma.
\]

These two identities imply $[A_g,I]=0$ without supposing that the coupled metric is vacuum Einstein. The full commutator in the common normalization is

\[
 [J,D_2](\gamma,\psi)=
 \begin{pmatrix}
 -g\sum_AV_A\psi^A\\
 (-V_A\operatorname{tr}_g\gamma)_A
 \end{pmatrix}.
 \tag{16}
\]

Indeed $2(IH_A-H_A)=-g\operatorname{tr}H_A=-gV_A$, and the lower mixed block gives $2H_A:I\gamma-2H_A:\gamma=-V_A\operatorname{tr}\gamma$. Thus under the hypothesis

\[
 V_A(\Phi(x))=0\quad\text{for every }A,x,
 \qquad [J,D_2]=0.
 \tag{17}
\]

The hypothesis implies $\Box\Phi^A=0$ and $V_{AB}d\Phi^B=0$, but does not imply $V_{AB}=0$. The latter identity merely says that the Hessian annihilates tangent directions to the background image. Nonconstant massless backgrounds and critical manifolds remain included.

\section{Cauchy reduction, charges and domains}

Choose a smooth compact Cauchy surface and a Gaussian collar $g=-dt^2+h_t$. Parallel transport the one-form and tensor bundles along its normal geodesics, use ordinary transport on the scalar factor, and conjugate by the scalar half-density ratio. This produces operators

\[
 D_i^{\mathrm{red}}=\partial_t^2+a_i(t)
\]

on fixed bundles over $\Sigma$ with fixed density $d\mathrm{vol}_{h_0}$. Their spatial principal symbols are $|\eta|_{h_t}^2\,1$. The conjugation is block diagonal between metric and scalar factors and carries $J$ to the time-independent $I_0\oplus1_N$.

Take the positive reference fibre metric on $E_2|_\Sigma$ to be four times the full Euclidean tensor contraction plus four times the Euclidean scalar contraction. In the split $(\gamma_{tt},\gamma_{t\Sigma},\gamma_{\Sigma\Sigma},\psi)$, put

\[
 \tau=\operatorname{diag}(+1,-1,+1,+1_N),\quad F=\tau J.
 \tag{18}
\]

Then the untrace-reversed Lorentzian bundle form is represented by $\tau$ and the physical form (8) by $F$. Both $I_0$ and $\tau$ are self-adjoint unitary involutions for the positive reference metric, and they commute. For example,

\[
 I_0\gamma=\gamma-\tfrac12g_0\langle g_0,\gamma\rangle_{\mathrm{Eucl}},
 \qquad \langle g_0,g_0\rangle_{\mathrm{Eucl}}=4,
\]

so $I_0$ is precisely a Euclidean reflection. It is not merely self-adjoint in an indefinite metric. Hence $F=F^*=F^{-1}$ as well.

For reduced Cauchy data $f=(u|_\Sigma,-i\partial_tu|_\Sigma)$ the charges, represented in the positive reference Hilbert pairing, are

\[
 \widetilde q=\begin{pmatrix}0&1\\1&0\end{pmatrix},\qquad
 q=\begin{pmatrix}0&\tau\\\tau&0\end{pmatrix},\qquad
 q_J=\begin{pmatrix}0&F\\F&0\end{pmatrix}=qJ_\Sigma.
 \tag{19}
\]

The charge current follows by integration by parts of the rough wave term. Zero-order Hermitian potentials cancel in the Green identity and add no boundary current. Therefore (19) is exact, not a principal-symbol approximation. The factor four in (8) has been included in the reference metric, not silently removed from the scalar or gravitational sector separately.

Let

\[
 \mathcal H_i^s=H^s(\Sigma;E_i)\oplus H^{s-1}(\Sigma;E_i).
\]

All statements of operator boundedness use these graded spaces. The forms $q_i$ and $q_J$ pair $\mathcal H^s$ continuously with $\mathcal H^{1-s}$. Quadratic positivity is asserted on smooth data, or on $\mathcal H^s$ for $s\ge1/2$, never on all Sobolev exponents.

With $U_i$ the Cauchy solution maps, define

\[
 K_\Sigma=\rho_2KU_1,\qquad
 K_\Sigma^\dagger=\rho_1K^\dagger U_2.
 \tag{20}
\]

Green intertwining and the Green identity show that the second operator is exactly the charged adjoint:

\[
 K_\Sigma^*q_J=q_1K_\Sigma^\dagger.
 \tag{21}
\]

They have bounded maps $K_\Sigma:\mathcal H_1^{s+1}\to\mathcal H_2^s$ and $K_\Sigma^\dagger:\mathcal H_2^s\to\mathcal H_1^{s-1}$. Normal derivatives of their first-order expressions are reduced using the equations. Finally,

\[
 K_\Sigma^\dagger K_\Sigma
 =\rho_1K^\dagger KU_1=\rho_1D_1U_1=0.
 \tag{22}
\]

In particular the gauge image is charge-isotropic and is contained in the constraint kernel. This step concerns solution Cauchy data; it is not the false spacetime identity $K^\dagger K=0$.

\section{Compatible microlocal projections}

Under (17), $D_2$ is formally self-adjoint for both the physical form $F$ and the untrace-reversed form $\tau$. Apply the general finite-rank Hermitian-bundle construction in G\'erard~\cite[Sections 3.1--3.2]{Gerard2024Compact} to the reduced operator using $\tau$. Compactness provides the uniform collar bounds used there. It gives a pseudodifferential Riccati factor $b\in\Psi^1$ with scalar positive principal root and

\[
 i\partial_tb-b^2+a_2\in\Psi^{-\infty},\qquad
 b+b^*\ge c\langle D\rangle,
 \tag{23}
\]

after a sufficiently large low-frequency cutoff. The adjoint $*$ in this inequality is the positive reference adjoint. The indefinite adjoint is $b^\star=\tau b^*\tau$.

The asymptotic Riccati recursion uniquely determines the symbol of $b$ modulo smoothing: at each descending order its unknown coefficient is multiplied by twice the nonzero scalar root. Since $[J,a_2]=0$ and $\partial_tJ=0$, the factor $JbJ$ has the same recursion and principal root. Thus $JbJ-b\in\Psi^{-\infty}$. Replace

\[
 b\longmapsto b_J=\tfrac12(b+JbJ).
 \tag{24}
\]

This preserves (23) and gives $[b_J,J]=0$ exactly. Its coercivity follows from

\[
 b_J+b_J^*=\tfrac12[(b+b^*)+J(b+b^*)J].
 \tag{25}
\]

For the Sobolev coercive norm one uses that $J$ is a smooth bounded invertible multiplier; alternatively choose the reference elliptic norm to commute with $J$. This avoids treating a pointwise unitary as automatically isometric for an arbitrarily chosen Sobolev norm.

Put $b^+=b_J$ and $b^-=-b_J^\star$. The difference $b^+-b^-=b_J+b_J^\star$ is an isomorphism $H^s\to H^{s-1}$: its positive-reference real part is the sum of the coercive real part of $b_J$ and its $\tau$-conjugate. Define the graph isomorphism

\[
 \mathscr F(v^+,v^-)=(v^++v^-,b^+v^++b^-v^-),
 \qquad
 c_2^\pm=\mathscr F\pi^\pm\mathscr F^{-1}.
 \tag{26}
\]

These are exact complementary projections on every $\mathcal H_2^s$. Their graph spaces are $q$-orthogonal because $b^-=-b^{+\star}$. The commutation with $J$ implies

\[
 q_Jc_2^\pm=c_2^{\pm *}q_J,\qquad [c_2^\pm,J_\Sigma]=0.
 \tag{27}
\]

Their graph forms also give directly

\[
 \langle(v,b^+v),\widetilde q(v,b^+v)\rangle
 =\langle v,(b_J+b_J^*)v\rangle\ge0,
\]

whereas the corresponding expression with $b^-$ is

\[
 -\langle v,\tau(b_J+b_J^*)\tau v\rangle\le0.
\]

Thus

\[
 \pm c_2^{\pm *}\widetilde q c_2^\pm\ge0.
 \tag{28}
\]

The Riccati construction yields $\mathrm{WF}'(U_2c_2^\pm)\subset\mathcal N^\pm\times T^*\Sigma$ (with the usual harmless noncharacteristic set in a local parametrix statement). Averaging by a smoothing term does not change this relation. Construct $c_1^\pm$ similarly for $D_1$; no trace-reversal symmetry is required for that bundle.

The exact differential intertwiner (15) implies the gauge compatibility

\[
 c_2^\pm K_\Sigma-K_\Sigma c_1^\pm\in\Psi^{-\infty}.
 \tag{29}
\]

For example, $KU_1c_1^+f$ solves $D_2u=0$ and has only positive-energy wavefront. Its negative-frequency component therefore has wavefront in the empty intersection $\mathcal N^+\cap\mathcal N^-$. The standard pseudodifferential Cauchy construction turns this into a smooth kernel on compact $\Sigma\times\Sigma$, exactly as in G\'erard~\cite[Lemma 3.2]{Gerard2024Compact}. This establishes smoothing, not merely a gain of one derivative.

\section{Positivity on the metric gauge slice}

Use the metric-only readout

\[
 l_N(\gamma,\psi)=(\operatorname{tr}_{g_0}\gamma,2\gamma_{t\Sigma}),
 \qquad Rf=(l_N\pi_0c_2^+f,l_N\pi_0c_2^-f).
 \tag{30}
\]

Its fibre kernel is independent of the scalar component. If $l_Nu_0=0$, then $I_0\gamma_0=\gamma_0$ and $\tau\gamma_0=\gamma_0$. Consequently $Fu_0=u_0$, including the scalar part. Since $F=F^*$, every datum $u=(u_0,u_1)$ satisfying this condition obeys

\[
 \langle u,q_Ju\rangle
 =\langle u_0,Fu_1\rangle+\langle u_1,Fu_0\rangle
 =\langle u_0,u_1\rangle+\langle u_1,u_0\rangle
 =\langle u,\widetilde qu\rangle.
 \tag{31}
\]

No condition on $u_1$ is needed. This is the exact charge identity, which is stronger and more relevant than merely checking that the field fibre metric is positive on a trace-free spatial slice.

For $f\in\ker R$, set $u^\pm=c_2^\pm f$. Each satisfies the hypothesis of (31). Equations (27), the projector identity and (28) give

\[
 \boxed{\ 
 \pm\langle f,q_Jc_2^\pm f\rangle
 =\pm\langle c_2^\pm f,q_Jc_2^\pm f\rangle
 =\pm\langle c_2^\pm f,\widetilde q c_2^\pm f\rangle\ge0,
 \quad f\in\ker R.
 \ }
 \tag{32}
\]

Thus the coupled positive-slice step is proved with a common normalization and exact Cauchy current. It neither assumes positivity of the full indefinite gauge bundle nor infers positivity from quotient nondegeneracy. The global gauge correction and spacetime lift are established in the following sections.


\setcounter{equation}{32}
\section{Complete gauge fixing and the Fredholm obstruction}\label{sec:fredholm}
We continue with the common normalization and reduced data above.
Write $\mathcal G^s=\Ran(K_\Sigma:\mathcal H_1^{s+1}\to\mathcal H_2^s)$ and $\mathcal C^s=\Ker K_\Sigma^\dagger$; omission of the superscript denotes smooth data.
Set $Y^s=H^s(\Sigma;E_1)\oplus H^s(\Sigma;E_1)$.
The readout $R:\mathcal H_2^s\to Y^s$ is surjective: it is the composition of the graph-coordinate isomorphism with $l_N\oplus l_N$, and $l_N$ has a smooth fibrewise right inverse.
For example, its prescribed trace can be supplied by a multiple of $g_0$ and its mixed component can be specified independently.
Surjectivity is an exact statement, including low frequencies.

\subsection{The exact boundary matrix}
Let $r=\tfrac12h_0^{-1}\dot h_0$ and $H=\tr r$.
For reduced one-form data $w=(w_t,w_\Sigma)$ the original covector is $a(t)W(t)w(t)$, where $a'/a=-H/2$ and spatial parallel transport satisfies $W'=rW$ at the reference slice.
The Gaussian Christoffel symbols are $\Gamma^j_{ti}=r_i{}^j$ and $\Gamma^t_{ij}=r_{ij}$.
Direct substitution gives
\begin{align*}
\tr_g(I d_s\xi)|_\Sigma
 &=\partial_t w_t+\tfrac H2w_t-\operatorname{div}_\Sigma w_\Sigma,\\
2(d_s\xi)_{t\Sigma}|_\Sigma
 &=\partial_t w_\Sigma-\tfrac H2w_\Sigma-rw_\Sigma+d_\Sigma w_t.
\end{align*}
Consequently, with $\delta_\Sigma=-\operatorname{div}_\Sigma$ and adjoints taken in the positive Riemannian product,
\begin{equation}\label{eq:B}
B=\begin{pmatrix}H/2&\delta_\Sigma\\d_\Sigma&-H/2-r\end{pmatrix},\qquad B^*=B,
\end{equation}
and
\begin{equation}\label{eq:readoutK}
l_N\pi_0K_\Sigma c_1^\pm
=\tfrac1{\sqrt2}(i b_1^\pm+B)\pi_0c_1^\pm.
\end{equation}
The coefficient of $r$ in~\eqref{eq:B} follows from the displayed differentiation.
In the primary rendering~\cite[Equations (4.10)--(4.11)]{Gerard2024Compact}, the derivative formula agrees with this coefficient whereas the following matrix displays $-r/2$.
We use~\eqref{eq:B}, with $r=\tfrac12h^{-1}\dot h$ fixed as above.

The scalar row of $K$ is removed exactly by $l_N$.
Using the smoothing intertwining defect and writing $L_1f=(\pi_0c_1^+f,\pi_0c_1^-f)$, we obtain
\begin{equation}\label{eq:RK}
A:=RK_\Sigma
=\tfrac1{\sqrt2}\begin{pmatrix}A^+&0\\0&A^-\end{pmatrix}L_1+S,
\quad A^\pm=i b_1^\pm+B,\quad S\in\Psi^{-\infty}.
\end{equation}
Thus all coupled lower-order dependence is included in $b_1^\pm$ and $S$; no unaccounted order-zero error is discarded.

\begin{proposition}\label{prop:fredholm}
For every $s\in\mathbb R$, $A:\mathcal H_1^{s+1}\to Y^s$ is Fredholm of index zero.
Its kernel and distributional cokernel are finite dimensional, smooth and independent of $s$.
The gauge range $\mathcal G^s$ is closed.
The induced charge on $\mathcal C/\mathcal G$ is nondegenerate.
\end{proposition}
\begin{proof}
Each $A^\pm:H^s\to H^{s-1}$ is elliptic: its principal symbol is $\pm i|\eta|+\sigma_1(B)$ and $\sigma_1(B)$ is Hermitian.
At $s=1/2$, the imaginary part of $(u,A^\pm u)$ is $(u,\operatorname{Re}b_1^\pm u)$, with a fixed sign and a coercive $H^{1/2}$ bound.
The forms rotated by $\mp i$ are coercive bounded forms on $H^{1/2}$.
Lax--Milgram gives injectivity and surjectivity $H^{1/2}\to H^{-1/2}$.
The adjoint estimate has the opposite imaginary sign and excludes a cokernel as well.
Elliptic regularity gives the isomorphism for every $s$.
Because $L_1:\mathcal H_1^{s+1}\to H^{s+1}\oplus H^{s+1}$ is an isomorphism, the diagonal term in~\eqref{eq:RK} is invertible onto $Y^s$.
Its smoothing perturbation is Fredholm of index zero, with smooth kernel and cokernel.
The perturbation can create both, so invertibility of $A$ is not assumed.

Choose a closed complement $X_0$ to $\Ker A$.
The restriction $A:X_0\to\Ran A$ is a bounded isomorphism.
If $K_\Sigma x_n$ converges for $x_n\in X_0$, then $Ax_n=RK_\Sigma x_n$ converges, hence $x_n$ converges.
Thus $K_\Sigma X_0$ is closed; adding the finite-dimensional smooth space $K_\Sigma\Ker A$ proves closedness of $\mathcal G^s$.

At $s=1/2$, the charge is a strong nondegenerate form: its off-diagonal entries identify $H^{1/2}$ and $H^{-1/2}$ with the appropriate dual spaces.
The charged-adjoint identity gives $\mathcal C^{1/2}=(\mathcal G^{1/2})^{q_J}$.
Closedness then implies $(\mathcal C^{1/2})^{q_J}=\mathcal G^{1/2}$, proving nondegeneracy of the quotient.
This also gives the assertion for smooth data.
Indeed, after Sobolev conjugation, the principal symbol of $K_\Sigma$ is injective because that of $RK_\Sigma$ is invertible.
Its positive normal operator is elliptic; its generalized inverse and range projection preserve smooth sections.
The corresponding constraint projection, obtained using the charge isomorphism, also preserves smooth sections.
Smooth constrained data are therefore dense in $\mathcal C^{1/2}$, and a smooth gauge-range element has a smooth gauge preimage.
\end{proof}

\section{Positive covariances, including the singular case}\label{sec:covariances}
We now construct a bounded gauge correction on each graded Sobolev space.
Here $Q=q_J$ denotes the physical Cauchy charge as a matrix in the positive auxiliary pairing.
For a field-data operator $T$, put $T^\dagger=Q^{-1}T^*Q$; for a readout map into the positive target, put $R^\dagger=Q^{-1}R^*$.
The distinction from $K_\Sigma^\dagger=q_1^{-1}K_\Sigma^*Q$ is essential.

\subsection{The regular case}
If $A$ is invertible, define
\[
T=1-K_\Sigma A^{-1}R.
\]
Then $RT=0$, $T^2=T$, $TK_\Sigma=0$, and $T$ preserves $\mathcal C$ because $K_\Sigma^\dagger K_\Sigma=0$.
For $f\in\mathcal C$, $Tf-f\in\mathcal G$.
The positive slice proved above gives positive Hermitian covariances
\[
\lambda_\Sigma^\pm=\pm T^*Q c_2^\pm T.
\]
They annihilate $\mathcal G$ in both arguments, and their difference is $Q$ on $\mathcal C\times\mathcal C$ by charge orthogonality to $\mathcal G$.
Moreover $c_2^\pm Tc_2^\mp$ is smoothing: $R$ separates the target frequencies, $A^{-1}$ is diagonal modulo smoothing in graph coordinates, and $K_\Sigma$ intertwines the frequency splittings modulo smoothing.

\subsection{The full kernel and cokernel construction}
Take a positive auxiliary $L^2$ basis $u_1,\ldots,u_n$ of $\Ker A^*$ and set
\[
v_j=R^\dagger u_j=Q^{-1}R^*u_j.
\]
Elliptic regularity makes all these vectors smooth, and surjectivity of $R$ makes the $v_j$ independent.
The exact charged-adjoint identity
\[
0=A^*u_j=K_\Sigma^*R^*u_j=q_1K_\Sigma^\dagger v_j
\]
places every $v_j$ in $\mathcal C$.
There is no assertion that $R^\dagger u_j$ has a vanishing scalar component.
The solvable-readout subspace is
\[
\mathcal H_{\mathrm{reg}}^s
=\{f:Rf\in\Ran A\}
=\{f:q_J(v_j,f)=0\text{ for every }j\}.
\]

Choose the basis so that $v_1,\ldots,v_p$ represent a complement of $V\cap\mathcal G$ in $V=\operatorname{span}(v_j)$ and $v_{p+1},\ldots,v_n\in\mathcal G$.
Nondegeneracy of $\mathcal C/\mathcal G$ and smooth density supply smooth $w_1,\ldots,w_p\in\mathcal C$ with an invertible pairing matrix against the first $p$ vectors.
Their pairings against the remaining vectors vanish.
Ambient charge nondegeneracy completes them by smooth $w_{p+1},\ldots,w_n$ to a full invertible matrix.
The block triangular inverse normalizes the dual family while keeping the first $p$ vectors in $\mathcal C$:
\begin{equation}\label{eq:dual}
q_J(v_j,w_i)=\delta_{ji},\qquad
w_i\in\mathcal C\ (i\le p),\quad v_j\in\mathcal C\ (j\le n).
\end{equation}
All dual pairings in~\eqref{eq:dual} use the physical charge.

Define smooth finite-rank projections
\[
\pi f=\sum_{j=1}^n w_jq_J(v_j,f),\qquad
\widetilde\pi f=\sum_{j=1}^p w_jq_J(v_j,f).
\]
Then $\Ker\pi=\mathcal H_{\mathrm{reg}}$, $\pi K_\Sigma=0$, and $\pi$ preserves $\mathcal C$.
On $\mathcal C$, $\pi=\widetilde\pi$, since the omitted $v_j$ belong to $\mathcal G$.
In addition,
\begin{equation}\label{eq:piannihilation}
\widetilde\pi K_\Sigma=\widetilde\pi^\dagger K_\Sigma=0.
\end{equation}
The second identity uses $w_j\in\mathcal C$ for $j\le p$.
The full $\pi^\dagger$ need only preserve $\mathcal G$ and is not asserted to annihilate it.

Let $\pi_1$ be a smooth finite-rank projection onto $\Ker A$ and let $G_A$ be the generalized inverse with range $\Ker\pi_1$, extended to all target data through a smooth cokernel projection.
Thus $G_AA=1-\pi_1$.
Choose a smooth finite-rank projection $\pi_2$ onto $K_\Sigma\Ker A=\mathcal G\cap\Ker R$.
Set
\begin{equation}\label{eq:Tsingular}
T_{\mathrm{reg}}=1-K_\Sigma G_AR\quad\hbox{on }\mathcal H_{\mathrm{reg}},
\qquad T=(1-\pi_2)T_{\mathrm{reg}}(1-\pi).
\end{equation}
These maps have the required boundedness on every $\mathcal H_2^s$.
Direct calculation gives $RT=0$, and $T$ is a projection onto $\Ker R\cap\Ker\pi_2$.
It preserves $\mathcal C$.
Since $T_{\mathrm{reg}}K_\Sigma=K_\Sigma\pi_1$, one has
\[
TK_\Sigma=(1-\pi_2)K_\Sigma\pi_1=0.
\]
The $\pi_2$ factor removes residual gauge fields from $\Ker A$ and cannot be omitted.
For constrained data,
\begin{equation}\label{eq:Tquotient}
Tf=(1-\pi)f\pmod{\mathcal G}
=(1-\widetilde\pi)f\pmod{\mathcal G}.
\end{equation}
All extra projections are smooth finite rank, so $c_2^\pm Tc_2^\mp$ remains smoothing.

Define the finite-rank Hermitian smoothing operator
\begin{equation}\label{eq:nu}
\nu=Q-(1-\widetilde\pi)^*Q(1-\widetilde\pi)
=Q\widetilde\pi+\widetilde\pi^*Q-\widetilde\pi^*Q\widetilde\pi.
\end{equation}
Equations~\eqref{eq:piannihilation} imply $\nu K_\Sigma=0$.
The positive and negative spectral parts in the auxiliary positive $L^2$ space,
\[
\nu_+=\tfrac12(|\nu|+\nu),\qquad
\nu_-=\tfrac12(|\nu|-\nu),
\]
are positive smooth finite-rank forms with $\nu=\nu_+-\nu_-$.
Their nonzero eigenvectors lie in the smooth range of $\nu$.
Every gauge vector belongs to its zero eigenspace, so $\nu_\pm$ both annihilate $\mathcal G$.
This conclusion follows by spectral calculus and requires no sign assumption on the finite physical obstruction space.

\begin{proposition}\label{prop:covariances}
The covariances
\begin{equation}\label{eq:lambdas}
\lambda_\Sigma^\pm=\pm T^*Q c_2^\pm T+\nu_\pm
\end{equation}
are positive and Hermitian, annihilate $\mathcal G$ and satisfy
$\lambda_\Sigma^+-\lambda_\Sigma^-=Q$ on $\mathcal C\times\mathcal C$.
Their cross-frequency errors are smoothing.
\end{proposition}
\begin{proof}
The first term in~\eqref{eq:lambdas} is positive because $\Ran T\subset\Ker R$ and the exact slice identity applies.
Hermiticity follows from physical-charge self-adjointness of $c_2^\pm$.
Gauge annihilation follows from $TK_\Sigma=0$ and the spectral argument.
Finally~\eqref{eq:Tquotient}, charge orthogonality to $\mathcal G$, and~\eqref{eq:nu} give, on constrained data,
\[
\lambda_\Sigma^+-\lambda_\Sigma^-
=T^*QT+\nu
=(1-\widetilde\pi)^*Q(1-\widetilde\pi)+\nu=Q.
\]
The finite-rank terms are smoothing and the remaining frequency assertion was proved after~\eqref{eq:Tsingular}.
\end{proof}

The construction treats gauge stabilizers $K_\Sigma\xi=0$, nonzero residual gauge fields $K_\Sigma\Ker A$, and genuine physical obstruction classes $V/(V\cap\mathcal G)$.
It follows the singular method of~\cite[Theorem 5.2]{Gerard2024Compact}, with a single physical charge throughout the dual matrix and defect~\eqref{eq:nu}.
Changing to the untrace-reversed charge in either step would not establish the stated identities.

\section{Spacetime lift and real physical state}\label{sec:lift}
Let $G_2=G_{2,\mathrm{ret}}-G_{2,\mathrm{adv}}$, with the charge convention fixed consistently with the Green identity, and let $\rho_2$ be the reduced trace.
For $f,h\in C_c^\infty(M;E_2)$ define
\begin{equation}\label{eq:lift}
\Lambda^\pm(f,h)=\lambda_\Sigma^\pm(\rho_2G_2f,\rho_2G_2h).
\end{equation}
The Green operator maps compact tests continuously to smooth spacelike-compact solutions, with advanced and retarded support in the corresponding causal cones.
Their Cauchy traces are smooth; compactness of $\Sigma$ admits every smooth datum.
Only a Gaussian collar is used for the microlocal construction, and neither a global Gaussian foliation nor geodesic completeness is needed.
Smooth finite-rank Cauchy terms lift to smooth kernels, since evolution sends smooth data to smooth global solutions.

\subsection{The physical test quotient}
Let $\Ker_cK^\dagger$ denote compactly supported smooth tests satisfying the constraint.
Intertwining implies that $\rho_2G_2$ maps this space into $\mathcal C$.
Furthermore $G_2Ph=-KG_1K^\dagger h$, so $P C_c^\infty(M;E_2)$ maps into the gauge range.
Hence~\eqref{eq:lift} descends to
\begin{equation}\label{eq:testquotient}
\frac{\Ker_cK^\dagger}{P C_c^\infty(M;E_2)}
\simeq\frac{\mathcal C}{\mathcal G}.
\end{equation}
We give the support-sensitive proof of the isomorphism.

Choose a smooth temporal cutoff $\chi$ whose transition lies in a compact Cauchy slab, with constant values zero in the past and one in the future.
Reverse its sign if required by the propagator convention.
Given $D_2u=0$ and $K^\dagger u=0$, put
\[
f=D_2(\chi u),\qquad v=K^\dagger(\chi u).
\]
Both have compact support because the equations vanish outside the transition.
Then $K^\dagger f=D_1v$, and
\[
f_{\mathrm{phys}}=f-Kv\in\Ker_cK^\dagger,
\qquad G_2f_{\mathrm{phys}}=u-KG_1v.
\]
The cutoff Green identity proves surjectivity onto the physical quotient, without requiring surjectivity onto each fixed constrained representative.

For injectivity, suppose $f\in\Ker_cK^\dagger$ and $G_2f=Kw$ with $D_1w=0$.
Set $\eta=D_1(\chi w)$ so that $G_1\eta=w$.
Then $G_2(f-K\eta)=0$.
Exactness of the normally hyperbolic causal sequence supplies compact $h$ satisfying $f-K\eta=D_2h$.
Applying $K^\dagger$ yields
\[
D_1(\eta+K^\dagger h)=0.
\]
This homogeneous solution has compact spacetime support and therefore vanishes by Cauchy uniqueness.
Consequently $f=D_2h-KK^\dagger h=Ph$.
Conversely every such $Ph$ maps to a gauge solution as above.
This proves~\eqref{eq:testquotient} using only causal-sequence exactness, cutoff Green identities and uniqueness for normally hyperbolic equations.

\subsection{Microlocal regularity and the neutral real field}
Because $c_2^\pm$ is $Q$-self-adjoint and $c_2^\pm Tc_2^\mp$ is smoothing, the same cross-frequency assertion holds for $T^\dagger$.
Thus
\[
U_2T^\dagger c_2^\pm T
=U_2c_2^\pm T^\dagger T\quad\pmod{\text{smooth kernel}}.
\]
The factor on the right acting on Cauchy data is graded pseudodifferential.
The microlocal Cauchy-lift criterion~\cite[Proposition 11.1.1]{Gerard2019Microlocal}, used in~\cite[proof of Theorem 5.1]{Gerard2024Compact}, together with propagation of singularities, gives
\[
\WF'(\Lambda^\pm)\subset\mathcal C_{\mathrm{Had}}^\pm
\subset\mathcal N^\pm\times\mathcal N^\pm.
\]
Here $\mathcal C_{\mathrm{Had}}^\pm$ is the relation of cotangent vectors on the same null bicharacteristic with the indicated energy sign, and $\WF'$ reverses the second covector.
Spatial pseudodifferential kernels are conormal to the diagonal and evolution propagates that relation along null bicharacteristics.
The equation removes any noncharacteristic contribution in the local parametrix statement.
The finite-rank corrections introduce no new singularities.
We assert this Hadamard inclusion for the redundant field representatives and exact commutation relations on physical tests; no equality on every component of the redundant bundle is required.

Arbitrary smoothing choices need not preserve the real-field structure.
On reduced Cauchy data let $C(f_0,f_1)=(\overline f_0,-\overline f_1)$ and define $a^C(f,h)=\overline{a(Cf,Ch)}$.
The real differential gauge complex is preserved and $Q^C=-Q$.
Replace the covariances by
\[
\lambda_{\mathrm{real}}^\pm
=\tfrac12\bigl(\lambda_\Sigma^\pm+(\lambda_\Sigma^\mp)^C\bigr).
\]
This preserves positivity, gauge annihilation and the exact physical commutation relations.
Conjugation exchanges the two characteristic signs, so the Hadamard property is preserved.
The pair now satisfies $(\lambda_{\mathrm{real}}^+)^C=\lambda_{\mathrm{real}}^-$ and defines a quasifree state on the neutral real physical CCR algebra.
Explicitly, on the $C$-fixed real quotient set $\sigma=-iQ$ and $\mu=\tfrac12(\lambda_{\mathrm{real}}^++\lambda_{\mathrm{real}}^-)$.
The relation $Q^C=-Q$ makes $\sigma$ real and antisymmetric, while positivity of $\lambda_{\mathrm{real}}^\pm=\mu\pm i\sigma/2$ gives the quasifree uncertainty inequality.
With the Weyl convention $W(f)W(h)=e^{-i\sigma(f,h)/2}W(f+h)$, the state is $\omega(W(f))=e^{-\mu(f,f)/2}$.
Together with Proposition~\ref{prop:covariances} and~\eqref{eq:testquotient}, this proves Theorem~\ref{thm:main}.

\section{A nonconstant coupled background with a massive spectator}\label{sec:example}
The critical-potential hypotheses do not reduce the theorem to pure gravity with constant scalar spectators.
Let $\gamma_{S^3}$ be the unit round metric and $Y:S^3\hookrightarrow\mathbb R^4$ its standard embedding.
The identities
\[
\sum_A(Y^A)^2=1,\quad
\sum_A dY^A\otimes dY^A=\gamma_{S^3},\quad
\Delta_{S^3}Y^A=-3Y^A
\]
give an exact closed Friedmann--Lema\^itre--Robertson--Walker ansatz
\[
g=-dt^2+a(t)^2\gamma_{S^3},\qquad
X^A(t,\omega)=f(t)Y^A(\omega),\quad A=1,\ldots,4,
\]
with flat potential $V=0$.
Write $\mathscr H=\dot a/a$ and $u=\dot f$ to distinguish the cosmological expansion rate from the spatial mean curvature used earlier.
The scalar equations and Einstein equations reduce to
\begin{align*}
\dot a&=a\mathscr H,&\dot f&=u,&
\dot u&=-3\mathscr H u-3f/a^2,\\
\dot{\mathscr H}&=a^{-2}-\tfrac12(\rho+p),&
\rho&=\tfrac12(u^2+3f^2/a^2),&p&=\tfrac12(u^2-f^2/a^2),
\end{align*}
with constraint $3(\mathscr H^2+a^{-2})=\rho+\Lambda$ in the canonical normalization.
The mixed momentum flow vanishes because $\sum_A Y^AdY^A=0$.
Differentiating gives $\dot\rho+3\mathscr H(\rho+p)=0$, and therefore propagation of the Einstein constraint.

At initial data
\[
(a,f,u,\mathscr H,\Lambda)=(1,1,3,1,0)
\]
one finds $\rho=6$, $p=4$, $\dot{\mathscr H}=-4$ and $\ddot f=-12$, satisfying the constraint and evolution equations.
Smooth local ODE existence with $a>0$ gives an actual on-shell spacetime $I\times S^3$, with compact Cauchy surfaces and no completeness requirement.
The metric--scalar coupling is nonzero:
\[
\nabla_t\nabla_iX^A=(u-\mathscr H f)\partial_iY^A
=2\partial_iY^A
\]
at the initial slice.
Moreover $dX=uY\,dt+f\,dY$ has rank four wherever $uf\ne0$, since the temporal direction is radial and the three spatial directions span $Y^\perp$.
Rank four is a property of this example and is not a hypothesis of the state theorem.

Adjoin a fifth background scalar $Z=0$ and potential $V(X,Z)=\mu^2Z^2/2$ with $\mu\ne0$.
It contributes no background stress and satisfies $V_{,A}=0$ on the entire image, while $V_{,ZZ}=\mu^2\ne0$.
The same nonconstant coupled solution therefore illustrates both flat background directions and a nonzero transverse potential Hessian.


\section{Scope, provenance and computational verification}
The construction combines an inherited microlocal gauge-fixing method with the explicitly derived finite-multiplet compatibility identities.
Its essential coupled features are the full scalar Gram and potential-Hessian blocks, the critical trace-reversal commutator, the common metric--scalar normalization, and the physical-charge finite-obstruction repair.
The finite-rank correction changes the state in a smooth sector and supplies no preferred state-selection prescription.

An accompanying exact-algebra script verifies the common mixed-adjoint normalization, the abstract charged-adjoint identity, the finite-rank physical-charge expansion, and the closed-cosmology constraint and background example.
These checks concern the stated algebraic identities.
Gauge annihilation of the spectral corrections follows from the general spectral argument in Section~\ref{sec:covariances}.
The pseudodifferential estimates and spacetime wavefront conclusion are analytic arguments in the text, using the cited primary results; they are not certified by a finite matrix computation or by a proof assistant.

\begin{samepage}
\paragraph{Contribution and review disclosure.}
The inherited work was initially Fable 5-led; according to the operator's account, results from July 2026 onward were predominantly GPT-produced.
DeepSeek and multiple GPT models also participated in review and preparation.
The retained source contains an 18 July 2026 checkpoint; the targeted audit and staging occurred on 28 October 2026, and the present consolidation is dated 29 October 2026.
These dates do not establish a public disclosure date or equation-by-equation model attribution.
Human authors and affiliations are intentionally blank, and human expert review has not occurred.
The operator's uncertainty is recorded verbatim:
\begin{quote}
I genuinely do not know how authorship, intellectual credit, or ownership should be attributed for these results.
\end{quote}
\end{samepage}

\begingroup
\footnotesize
\bibliographystyle{plainurl}
\bibliography{references}
\endgroup
\end{document}
